\section{Kategorie i ich pierwsze przykłady}\label{sec:kategorie}

\subsection{Definicja kategorii}

\begin{definicja}[Kategoria]\label{def:kategoria}
\emph{Kategorią} $\mathcal C$ nazywamy dane obejmujące:
\begin{enumerate}
\item klasę obiektów $\Ob(\mathcal C)$;
\item dla każdej uporządkowanej pary obiektów $A,B$ kolekcję morfizmów $\Hom_{\mathcal C}(A,B)$; w tym skrypcie zakładamy, że jest ona zbiorem; każdy morfizm traktujemy wraz z jego dziedziną $A$ i przeciwdziedziną $B$;
\item dla każdej trójki obiektów $A,B,C$ operację składania
$\Hom_{\mathcal C}(B,C)\times\Hom_{\mathcal C}(A,B)\to\Hom_{\mathcal C}(A,C)$, $(g,f)\mapsto g\circ f$;
\item dla każdego obiektu $A$ wyróżniony morfizm $\id_A\in\Hom_{\mathcal C}(A,A)$.
\end{enumerate}
Dane te muszą spełniać dwa aksjomaty:
\begin{enumerate}[label=(\roman*)]
\item \emph{Łączność}: jeśli $f:A\to B$, $g:B\to C$ i $h:C\to D$, to $h\circ(g\circ f)=(h\circ g)\circ f$.
\item \emph{Jednostki}: dla każdego $f:A\to B$ zachodzi $f\circ\id_A=f=\id_B\circ f$.
\end{enumerate}
Kategorię nazywamy \emph{lokalnie małą}, jeśli dla każdej pary obiektów jej kolekcja morfizmów jest zbiorem; to właśnie konwencja przyjęta dalej. Kategoria jest \emph{mała}, jeśli jej obiekty i wszystkie morfizmy tworzą zbiory. W kategorii lokalnie małej o zbiorze obiektów zbiór wszystkich morfizmów otrzymujemy jako sumę rozłączną zbiorów $\Hom_{\mathcal C}(A,B)$ indeksowaną parami obiektów. Kategorię, która nie jest mała, nazywamy dużą.
\end{definicja}

Morfizmy nie muszą być funkcjami pomiędzy zbiorami. Aksjomaty wymagają tylko możliwości ich składania i istnienia identyczności. W zapisie $g\circ f$ najpierw wykonujemy $f$, a następnie $g$.

\begin{lemat}[Jednoznaczność identyczności]
Jeśli $e,e':A\to A$ są morfizmami spełniającymi aksjomat jednostki dla wszystkich morfizmów o odpowiednich dziedzinach i przeciwdziedzinach, to $e=e'$.
\end{lemat}
\begin{proof}
Z faktu, że $e$ jest jednostką lewostronną, wynika $e\circ e'=e'$. Z faktu, że $e'$ jest jednostką prawostronną, wynika $e\circ e'=e$. Zatem $e=e'$.
\end{proof}


\begin{uwaga}[Kwestie teoriomnogościowe]
Obiekty kategorii nie muszą tworzyć zbioru: na przykład ``wszystkie zbiory'' tworzą klasę właściwą. Przyjmujemy standardową konwencję, że rozpatrywane kategorie są lokalnie małe, a słowo \emph{mały diagram} oznacza diagram indeksowany małą kategorią. Gdy chcemy tworzyć kategorię wszystkich funktorów pomiędzy dużymi kategoriami, trzeba dodatkowo kontrolować wielkość kolekcji transformacji naturalnych. Przy głównych zastosowaniach rozróżnienie to zaznaczymy jawnie.
\end{uwaga}


\subsection{Wielkość kategorii}

\begin{definicja}[Uniwersum Grothendiecka]\label{cat:def:uniwersum}
Uniwersum Grothendiecka $\mathcal U$ to zbiór przechodni domknięty na tworzenie par, zbiorów potęgowych oraz sum rodzin $(x_i)_{i\in I}$ takich, że $I\in\mathcal U$ i $x_i\in\mathcal U$ dla każdego $i\in I$. Przyjmujemy ponadto $\mathbb N\in\mathcal U$. Przechodniość oznacza, że $x\in y\in\mathcal U$ implikuje $x\in\mathcal U$. Zbiory należące do $\mathcal U$ nazywamy $\mathcal U$-małymi.
\end{definicja}

\begin{uwaga}[Relatywność pojęcia małości]\label{cat:rem:uniwersa}
Jeżeli $X,Y\in\mathcal U$, to $X\times Y\in\mathcal U$ (pary uporządkowane tworzymy za pomocą par nieuporządkowanych, a zbiór wszystkich takich par jest podzbiorem odpowiedniego iterowanego zbioru potęgowego). Każdy podzbiór $X\times Y$ również należy do $\mathcal U$: należy on do $\mathcal P(X\times Y)$, a ta należy do $\mathcal U$. W szczególności zbiór funkcji $X\to Y$, utożsamianych z ich wykresami, jest $\mathcal U$-mały, gdyż jest podzbiorem $\mathcal P(X\times Y)$. Natomiast kolekcja wszystkich $\mathcal U$-małych zbiorów to samo $\mathcal U$, które nie należy do siebie na mocy aksjomatu ufundowania. Kategoria tych zbiorów jest więc lokalnie $\mathcal U$-mała, ale nie jest $\mathcal U$-mała. Może być mała względem większego uniwersum zawierającego $\mathcal U$.

Istnienia dowolnie dużych uniwersów tego rodzaju nie przyjmujemy jako konsekwencji standardowych aksjomatów ZFC. Jeżeli posługujemy się uniwersami, odpowiednie założenie teoriomnogościowe wskazujemy jawnie; w pozostałym tekście wystarczy konwencja klas i zbiorów z definicji~\ref{def:kategoria}.
\end{uwaga}

\subsection{Kategorie znane z innych dziedzin}

\begin{przyklad}[Zbiory]
W kategorii $\Set$ obiektami są zbiory, morfizmami są funkcje, a składanie i identyczności są zwykłymi operacjami na funkcjach. Jest to przykład kategorii lokalnie małej, ale nie małej: nie istnieje zbiór wszystkich zbiorów.
\end{przyklad}

\begin{przyklad}[Grupy i monoidy]
W $\Grp$ obiektami są grupy, a morfizmami homomorfizmy grup. W $\Mon$ obiektami są monoidy (zbiory z łącznym działaniem i elementem neutralnym), a morfizmami homomorfizmy zachowujące działanie i element neutralny. Złożenie homomorfizmów i identyczność są homomorfizmami, dlatego otrzymujemy kategorie.
\end{przyklad}

\begin{przyklad}[Przestrzenie liniowe]
Ustalmy ciało $k$. W kategorii $\Vect_k$ obiektami są przestrzenie wektorowe nad $k$, a morfizmami odwzorowania $k$-liniowe. Kategoria $\Vect_k^{\mathrm{fd}}$ to jej pełna podkategoria przestrzeni skończenie wymiarowych: oznacza to, że rozpatrujemy tylko takie obiekty, lecz wszystkie odwzorowania liniowe pomiędzy nimi. W kategorii $\Ab$ obiektami są grupy abelowe, a morfizmami homomorfizmy.
\end{przyklad}

\begin{przyklad}[Topologia i pierścienie]
W $\Top$ obiektami są przestrzenie topologiczne, a morfizmami odwzorowania ciągłe. W $\Ring$ rozważamy pierścienie przemienne z jedynką i homomorfizmy zachowujące $1$. Ważne jest określenie konwencji: bez podania, jakie pierścienie i jakie homomorfizmy rozważamy, nazwa kategorii nie wyznacza całej struktury.
\end{przyklad}

\begin{przyklad}[Moduły i rozmaitości różniczkowe]\label{cat:ex:moduly-rozmaitosci}
Dla pierścienia $R$ z jedynką oznaczmy przez $\cat{Mod}_R^{\mathrm{l}}$ kategorię lewych unitarnych $R$-modułów i homomorfizmów $R$-liniowych. Nie zakładamy przemienności $R$. Jeśli $f:M\to N$ i $g:N\to P$ są liniowe, to dla $r,s\in R$ oraz $x,y\in M$ zachodzi
\[
(g\circ f)(rx+sy)=g(rf(x)+sf(y))
=r(g\circ f)(x)+s(g\circ f)(y).
\]
Identyczność również jest liniowa, więc składanie funkcji wyznacza kategorię. Jest to przykład oparty na konstrukcjach z rozdziału 16 skryptu \emph{Algebra abstrakcyjna}.

W kategorii $\cat{Man}$ obiektami są gładkie rozmaitości bez brzegu (Hausdorffa, z przeliczalną bazą topologii), a morfizmami są gładkie odwzorowania. Złożenie jest gładkie, ponieważ w mapach lokalnych jest złożeniem gładkich funkcji między otwartymi podzbiorami przestrzeni euklidesowych; identyczność w tej samej mapie lokalnej jest identycznością euklidesową. Korzystamy tu z pojęć w rozdziale 1 skryptu \emph{Geometria i topologia różniczkowa}.
\end{przyklad}

\begin{definicja}[Kategoria konkretna]\label{cat:def:konkretna}
Kategoria konkretna nad $\Set$ to kategoria $\mathcal C$ wyposażona w przyporządkowanie każdemu obiektowi $A$ zbioru $U(A)$ i każdemu morfizmowi $f:A\to B$ funkcji $U(f):U(A)\to U(B)$, zachowujące identyczności i składanie oraz rozróżniające morfizmy o ustalonych dziedzinie i przeciwdziedzinie: $U(f)=U(g)$ implikuje $f=g$. W języku rozdziału~\ref{sec:funktory} jest to wierny funktor $U:\mathcal C\to\Set$.
\end{definicja}

\begin{przyklad}
Kategorie grup, modułów, przestrzeni topologicznych i rozmaitości są konkretne dzięki przyporządkowaniu zbioru bazowego oraz tej samej funkcji bazowej. Konkretność nie wymaga, aby każda funkcja między zbiorami bazowymi zachowywała strukturę: np. nie każda funkcja między grupami jest homomorfizmem.
\end{przyklad}

\begin{przyklad}[Kategoria dyskretna]
Dla zbioru $S$ zdefiniujmy kategorię $\mathcal D_S$ o obiektach $S$, w której $\Hom(x,x)=\{\id_x\}$, a $\Hom(x,y)=\varnothing$ dla $x\neq y$. Morfizmami są wyłącznie identyczności. Taka kategoria przechowuje informację o elementach $S$, ale nie o relacjach między nimi.
\end{przyklad}

\begin{przyklad}[Zbiory uporządkowane]\label{przyklad:poset}
Niech $(P,\leq)$ będzie zbiorem częściowo uporządkowanym. Obiektami kategorii $\mathcal P$ są elementy $P$, a $\Hom_{\mathcal P}(x,y)$ jest zbiorem jednoelementowym, gdy $x\leq y$, i pustym w przeciwnym przypadku. Zwrotność zapewnia identyczności, przechodniość zapewnia składanie, a łączność wynika z tego, że między ustalonymi obiektami istnieje co najwyżej jeden morfizm.

Dla przykładu zbiór $\{1,2,3,6\}$ uporządkowany przez podzielność ma morfizm $2\to 6$, ale nie ma morfizmu $6\to 2$. Nie należy interpretować tych strzałek jako funkcji pomiędzy liczbami. Każdy \emph{praporządek} (relacja zwrotna i przechodnia, niekoniecznie antysymetryczna) także definiuje kategorię.
\end{przyklad}

\begin{przyklad}[Monoid jako kategoria z jednym obiektem]\label{przyklad:monoid-kat}
Niech $(M,\cdot,1)$ będzie monoidem. Tworzymy kategorię z jednym obiektem $*$, przyjmując $\Hom(*,*)=M$. Składanie definiujemy przez $g\circ f=g\cdot f$, a identycznością jest $1$. Aksjomaty kategorii to dokładnie łączność i istnienie jednostki w monoidzie.

W szczególności grupa $G$ jest kategorią z jednym obiektem, której wszystkie morfizmy są odwracalne. Tym samym teoria kategorii zawiera teorię monoidów i grup jako szczególne przypadki.
\end{przyklad}

\begin{definicja}[Grupoid]\label{cat:def:grupoid}
Grupoid to kategoria, w której każdy morfizm $f:A\to B$ ma odwrotny morfizm $g:B\to A$, czyli $g\circ f=\id_A$ oraz $f\circ g=\id_B$. Taki morfizm $g$ jest jedyny: dla drugiej odwrotności $g'$ mamy $g=g\circ\id_B=g\circ(f\circ g')=(g\circ f)\circ g'=\id_A\circ g'=g'$.
\end{definicja}

\begin{przyklad}[Grupoidy]\label{cat:ex:grupoid}
Kategoria dyskretna jest grupoidem. Grupa rozpatrywana jako kategoria z jednym obiektem również jest grupoidem. Dla dowolnego zbioru $S$ możemy zbudować grupoid par: obiektami są elementy $S$, a jedynym morfizmem $a\to b$ jest para $(a,b)$. Definiujemy $(b,c)\circ(a,b)=(a,c)$ oraz $\id_a=(a,a)$. Oba nawiasowania złożenia trzech strzałek dają tę samą parę z początkiem pierwszej i końcem ostatniej strzałki; identyczności nie zmieniają tych końców. Odwrotnością $(a,b)$ jest $(b,a)$. Inny ważny przykład stanowi grupoid podstawowy przestrzeni topologicznej, którego strzałkami są klasy homotopii dróg o ustalonych końcach (por. rozdział 4 skryptu \emph{Topologia algebraiczna}); jego formalną konstrukcję odkładamy do części o zastosowaniach.
\end{przyklad}

\begin{przyklad}[Macierze i standardowe współrzędne]\label{przyklad:mat}
W kategorii $\cat{Mat}_k$ obiektami są liczby $0,1,2,\ldots$, a morfizmami $n\to m$ macierze o $m$ wierszach i $n$ kolumnach z wyrazami w $k$. Składaniem jest mnożenie macierzy w tej samej kolejności co składanie odwzorowań, a identycznością macierz jednostkowa $I_n$ (dla $n=0$ pusta macierz $0\times0$).

Przykładowo macierz $\begin{pmatrix}1&2\\0&1\\3&0\end{pmatrix}$ jest morfizmem $2\to3$. Kategoria ta koduje wszystkie skończenie wymiarowe przestrzenie liniowe \emph{z wyborem wzorcowego modelu} $k^n$.
\end{przyklad}

\begin{przyklad}[Relacje zamiast funkcji]
W $\Rel$ obiektami są zbiory, a morfizm $R:A\to B$ to relacja $R\subseteq A\times B$. Dla $R:A\to B$ i $S:B\to C$ definiujemy $S\circ R=\{(a,c):\text{istnieje }b\in B\text{ z }(a,b)\in R,(b,c)\in S\}$. Identycznością na $A$ jest przekątna $\Delta_A=\{(a,a):a\in A\}$. Aby wykazać łączność, niech dodatkowo $T:C\to D$ będzie relacją. Dla $a\in A$ i $d\in D$ warunek $(a,d)\in T\circ(S\circ R)$ jest równoważny istnieniu $b\in B$ i $c\in C$ spełniających
\[
(a,b)\in R,\qquad (b,c)\in S,\qquad (c,d)\in T.
\]
Dokładnie ten sam warunek opisuje przynależność do $(T\circ S)\circ R$. Ponadto $(a,b)\in R\circ\Delta_A$ wtedy i tylko wtedy, gdy istnieje $a'\in A$ takie, że $a=a'$ i $(a',b)\in R$, czyli gdy $(a,b)\in R$. Analogicznie $\Delta_B\circ R=R$.

Na przykład dla $A=B=C=\{1,2\}$ oraz $R=\{(1,1),(1,2)\}$, $S=\{(1,2),(2,2)\}$ otrzymujemy $S\circ R=\{(1,2)\}$. Kategoria $\Rel$ pokazuje, że kategorie nie muszą być zbudowane z funkcji.
\end{przyklad}


\begin{przyklad}[Ścieżki w grafie skierowanym]\label{przyklad:sciezki}
Rozważmy graf z wierzchołkami $a,b,c$ i krawędziami $u:a\to b$, $v:b\to c$. Obiektami kategorii ścieżek są wierzchołki; morfizmami są skończone skierowane ścieżki, w tym ścieżki długości $0$ (identyczności). W szczególności istnieje morfizm $v\circ u:a\to c$, chociaż w wyjściowym grafie nie było osobnej krawędzi $a\to c$. Składanie polega na konkatenacji skończonych ciągów krawędzi: najpierw zapisujemy krawędzie pierwszej ścieżki, potem drugiej. Dwa nawiasowania konkatenacji trzech ciągów dają dokładnie ten sam ciąg, a dopisanie pustego ciągu w odpowiednim wierzchołku nie zmienia ścieżki. Otrzymujemy przykład \emph{kategorii swobodnej na grafie}. Pokazuje on, że sam graf skierowany nie jest jeszcze kategorią: trzeba uwzględnić również wszystkie wymagane złożenia.
\end{przyklad}

\begin{przyklad}[Przestrzenie z wyróżnionym punktem]
W kategorii $\Top_*$ obiektami są pary $(X,x_0)$ z przestrzenią topologiczną $X$ i wyróżnionym punktem $x_0\in X$. Morfizm $f:(X,x_0)\to(Y,y_0)$ to odwzorowanie ciągłe spełniające $f(x_0)=y_0$. Przykładowo funkcja $t\mapsto t^2$ daje morfizm $([0,1],0)\to([0,1],0)$, ale nie daje morfizmu $([0,1],1)\to([0,1],0)$. Ten sam wzór funkcji może więc być morfizmem w jednej kategorii, a w innej nie.
\end{przyklad}

\subsection{Podkategorie i kategoria przeciwna}

\begin{definicja}[Podkategoria]
\emph{Podkategorią} $\mathcal D$ kategorii $\mathcal C$ nazywamy kategorię, której obiekty są obiektami $\mathcal C$, zbiory morfizmów $\Hom_{\mathcal D}(A,B)$ są podzbiorami $\Hom_{\mathcal C}(A,B)$, a identyczności i składanie są odziedziczone z $\mathcal C$. Jest \emph{pełna}, jeśli dla dowolnych $A,B\in\Ob(\mathcal D)$ zachodzi $\Hom_{\mathcal D}(A,B)=\Hom_{\mathcal C}(A,B)$.
\end{definicja}

\begin{przyklad}
Przestrzenie skończenie wymiarowe tworzą pełną podkategorię $\Vect_k$. Z kolei kategoria mająca te same obiekty co $\Vect_k$, ale jedynie liniowe izomorfizmy jako morfizmy, nie jest pełną podkategorią $\Vect_k$: pomija np. odwzorowanie zerowe $k\to k$.
\end{przyklad}

\begin{definicja}[Kategoria przeciwna]\label{def:op}
Dla kategorii $\mathcal C$ \emph{kategorię przeciwną} $\mathcal C^{\op}$ otrzymujemy, pozostawiając te same obiekty i odwracając kierunki morfizmów: $\Hom_{\mathcal C^{\op}}(A,B)=\Hom_{\mathcal C}(B,A)$. Piszemy $f^{\op}:B\to A$ na oznaczenie strzałki odpowiadającej $f:A\to B$ w $\mathcal C$. Jeżeli $f:A\to B$ oraz $g:B\to C$, definiujemy $f^{\op}\circ_{\mathcal C^{\op}}g^{\op}=(g\circ_{\mathcal C}f)^{\op}$ oraz $\id_A^{\op}$ jako identyczność w kategorii przeciwnej. Dla dodatkowego $h:C\to D$ mamy
\[
(f^{\op}\circ g^{\op})\circ h^{\op}
=(h\circ(g\circ f))^{\op}
=((h\circ g)\circ f)^{\op}
=f^{\op}\circ(g^{\op}\circ h^{\op}).
\]
Obie tożsamości jednostkowe wynikają z $(f\circ\id_A)^{\op}=f^{\op}=(\id_B\circ f)^{\op}$. Zatem kategoria przeciwna spełnia aksjomaty kategorii.
\end{definicja}

\begin{uwaga}
Dwukrotne odwrócenie strzałek daje $(\mathcal C^{\op})^{\op}=\mathcal C$. Jeśli kategoria pochodzi ze zbioru uporządkowanego, jej kategoria przeciwna odpowiada porządkowi odwrotnemu. Pojęcie \emph{dualności} polega na tym, że każde twierdzenie wynikające wyłącznie z aksjomatów kategorii ma wersję z odwróconymi wszystkimi strzałkami.
\end{uwaga}

